\begin{center}
{\heiti\zihao{3}地形与能量约束下的山区无人机运输、通信与资源分区}\par
\vspace{12pt}
{\heiti\zihao{4}摘要}\par
\end{center}
针对山区 \NumServices{} 个服务区、\NumBoxes{} 个不可拆货箱的运输任务，本文以 DEM 航高、动态剩余载荷和分阶段能量构成统一物理底座，将能力边界、资源调度、通信时序和独立配置连成递进模型。

针对 Q1，先扣除地形爬升、空载返程与返航储备，利用载荷—能耗单调性二分求安全载荷，再对满足质量、体积和能量约束的货箱子集进行集合划分动态规划。得到 \QOneTrips{} 架次、\QOneEnergy{} kWh、累计作业 \QOneTime{} s；精确性限于单点往返及“架次—能耗—累计时间”词典序目标。

针对 Q2，加入实体无人机、共享电池充电、硬截止与软延误，采用紧迫任务优先构造、可行合并和资源重排，保存 \CandidateCount{} 个有限候选。显式选择 \QTwoTrips{} 架次方案，能耗 \QTwoEnergy{} kWh，完工 \QTwoTime{} s，全部硬截止满足；相邻候选比较说明继续压缩架次须付出及时性代价，该选择不构成全局最优证明。

针对 Q3，按双向预算与 DEM 遮挡检查飞行、交接全过程，以指定单跳中继覆盖直连不足区间。在固定分组、路线、资源顺序和中继位置/高度下，用词典序 LP 优化运输开始和中继窗口，得到 \QThreeTrips{} 个运输架次、\RelayTrips{} 个中继架次，联合能耗 \JointEnergy{} kWh、完工 \JointTime{} s。1 s 与 0.5 s 回放均未发现中断，最小硬截止余量 \HardSlack{} s；最优性仅针对固定方案的时序问题。

针对 Q4，继承 Q3 的任务与时序，将服务区归并为 \AtomicGroups{} 个不可拆连通分量，以区间峰值核算资源，精确枚举两组和三组分区，缺口分别为 \GapTwo{} 与 \GapThree{}。该结果以组内同型资源可重编号、中继完整任务可按组复制为前提。

输入摘要核验、逐箱逐资源回放与边界反例构成验证链。但加密通信最小余量仅 \MarginHalf{} dB，固定运输计划的返航电量仅容许能耗统一上浮约 \QThreeEnergyTolerance\%。结果支持官方确定性场景中的可复现候选，不能替代连续通信证明或实飞验证。

\vspace{8pt}
\noindent\textbf{关键词：}山区无人机运输；动态载荷能耗；集合划分；共享电池调度；通信中继；独立资源配置
